% Generated by presentation/scripts/make_figures.py from presentation/data/*.csv. Needs \usepackage{booktabs}.

\begin{tabular}{p{6.4cm}lrrr}
\toprule
Data & Stage & Episodes & Examples & Per shape \\
\midrule
human T demonstrations (original dataset, not new) & all stages & 206 & 24,208 & T only \\
teacher demonstrations, 5 shapes (noisy-state teacher, every step) & 5 shapes & 500 & 154,170 & -- \\
teacher demonstrations for the 5 added letters, executed-push labels & 10 shapes & 500 & -- & 100 \\
teacher demonstrations with intended 16-step plans, 10 shapes & 10 shapes & 2,500 & 357,660 & 250 \\
teacher demonstrations with intended plans for the 17 new letters & 27 shapes & 2,550 & 312,144 & 150 \\
teacher demonstrations, all 27 shapes & 27 shapes & 5,050 & 669,804 & -- \\
correction scenes per shape per round & 10 and 27 shapes & -- & -- & 30 \\
correction scenes per new letter per round, from round 14 & 27 shapes & -- & -- & 60 \\
correction examples after round 50 (26-letter 96 px line) & 27 shapes & -- & 2,511,488 & -- \\
all training examples in round 50 (26-letter 96 px line) & 27 shapes & -- & 3,205,500 & -- \\
training samples per shape per round, 10 shapes, 20k steps & 10 shapes & -- & -- & 512k \\
training samples per shape per round, 27 shapes, 20k steps & 27 shapes & -- & -- & 190k \\
training samples per shape per round, 27 shapes, 54k steps & 27 shapes & -- & -- & 512k \\
\bottomrule
\end{tabular}
% presentation/data/dataset_sizes.csv (sources per row): episodes = scenes, examples = training samples (one per recorded step); per shape = episodes, scenes per round or samples per round as the row says.
